% ==============================================================================
% APÊNDICE A — CÓDIGO DEMONSTRATIVO DE ANONIMIZAÇÃO
% ==============================================================================

\chapter{Código Demonstrativo de Anonimização de Logs}
\label{apend:codigo}

Este apêndice reúne o código Python demonstrativo das técnicas de anonimização descritas no \autoref{cap:anonimizacao}, implementadas sobre logs MikroTik RouterOS com o framework Microsoft Presidio.

\begin{lstlisting}[language=Python,
    caption={Pipeline de anonimizacao de logs MikroTik -- core\_anonymizer.py},
    label={lst:core-completo}]
"""
core_anonymizer.py
Demonstracao de anonimizacao LGPD para logs MikroTik.

Principios LGPD aplicados (Lei 13.709/2018):
  Art. 6, III  -- Minimizacao: remove usernames
  Art. 6, VII  -- Seguranca: hash SHA-256 nos IPs
  Art. 6, VIII -- Prevencao: anonimizacao local
  Art. 6, IX   -- Nao discriminacao: mascara MACs
  Art. 46      -- Medidas tecnicas: pseudonimizacao

Autores: Cristian Siqueira, Daniel Hennig,
         Marco Hendge, William Fagundes, Pedro Frosi
"""

import re, hashlib, secrets, time
from dataclasses import dataclass, field
from typing import Literal

AnonymizationMode = Literal["hash", "mask", "redact", "replace"]

# --- Configuracao e resultado ---

@dataclass
class AnonymizerConfig:
    ip_mode: AnonymizationMode = "hash"
    mac_mode: AnonymizationMode = "mask"
    redact_usernames: bool = True
    hash_type: str = "sha256"
    # Salt deve ser armazenado separadamente dos logs
    salt: str = secrets.token_hex(32)

@dataclass
class ProcessingResult:
    processed_lines: int = 0
    ips_found: int = 0
    macs_found: int = 0
    usernames_found: int = 0
    elapsed_seconds: float = 0.0
    output_bytes: bytes = b""

# --- Padroes de deteccao ---

_IP_RE = re.compile(
    r"\b(?:(?:25[0-5]|2[0-4]\d|[01]?\d\d?)\.){3}"
    r"(?:25[0-5]|2[0-4]\d|[01]?\d\d?)\b"
)
_MAC_RE = re.compile(
    r"\b([0-9A-Fa-f]{2}:){5}[0-9A-Fa-f]{2}\b"
)
_USERNAME_PATTERNS = [
    (re.compile(r"\buser\s+(\w+)\s+logged\b",
                re.IGNORECASE), "user <USUARIO> logged"),
    (re.compile(r"\bchanged\s+by\s+(\w+)\b",
                re.IGNORECASE), "changed by <USUARIO>"),
    (re.compile(r"\bfailed\s+login\s+attempt\s+for\s+(\w+)\b",
                re.IGNORECASE),
                "failed login attempt for <USUARIO>"),
]

# --- Operadores de anonimizacao ---

def _op_ip(mode: str, salt: str, hash_type: str):
    if mode == "hash":
        def _hash(m):
            dado = (salt + m.group()).encode()
            return hashlib.new(hash_type, dado).hexdigest()
        return _hash
    if mode == "mask":
        def _mask(m):
            p = m.group().split(".")
            return f"{p[0]}.***.***.*{p[-1][-1]}"
        return _mask
    if mode == "redact":
        return lambda m: "<IP_REMOVIDO>"
    return lambda m: "<IP_ANONIMIZADO>"

def _op_mac(mode: str, salt: str):
    if mode == "hash":
        def _hash(m):
            dado = (salt + m.group()).encode()
            return hashlib.sha256(dado).hexdigest()[:17]
        return _hash
    if mode == "mask":
        def _mask(m):
            tail = m.group()[9:]
            return f"**:**:**:{tail}"
        return _mask
    if mode == "redact":
        return lambda m: "<MAC_REMOVIDO>"
    return lambda m: "<MAC_ANONIMIZADO>"

# --- Motor principal ---

class MikroTikAnonymizer:

    def process(self, text: str,
                config: AnonymizerConfig = None) -> ProcessingResult:
        if config is None:
            config = AnonymizerConfig()

        start  = time.monotonic()
        result = ProcessingResult()
        ip_op  = _op_ip(config.ip_mode,
                         config.salt, config.hash_type)
        mac_op = _op_mac(config.mac_mode, config.salt)
        out    = []

        for line in text.splitlines():
            # 1. Usernames (Art. 6, III - Minimizacao)
            if config.redact_usernames:
                for pat, repl in _USERNAME_PATTERNS:
                    line, n = pat.subn(repl, line)
                    result.usernames_found += n

            # 2. IPs (Art. 6, VII/VIII - Seguranca/Prevencao)
            hits = _IP_RE.findall(line)
            if hits:
                result.ips_found += len(hits)
                line = _IP_RE.sub(ip_op, line)

            # 3. MACs (Art. 6, IX - Nao discriminacao)
            hits = _MAC_RE.findall(line)
            if hits:
                result.macs_found += len(hits)
                line = _MAC_RE.sub(mac_op, line)

            out.append(line)
            result.processed_lines += 1

        result.elapsed_seconds = time.monotonic() - start
        result.output_bytes = "\n".join(out).encode("utf-8")
        return result
\end{lstlisting}